% ============================================================
% FUENTES: "Notas Biomate 1.pdf" (español) — NO está en el draft.
% ============================================================
\section{Aplicación: redes neuronales y función de activación sigmoidal}

Como aplicación biológica y computacional directa de la bifurcación pitchfork, consideremos el modelo dinámico de una neurona artificial o ensamble neuronal auto-recurrente con memoria.

\subsection{Motivación: redes recurrentes (RNN) y gradiente desvaneciente}

En el aprendizaje profundo y en los modelos continuos de neurociencia computacional, una de las funciones de activación más emblemáticas para mitigar el desvanecimiento del gradiente en redes recurrentes (\emph{Recurrent Neural Networks}, RNN) es la \textbf{tangente hiperbólica}:
\begin{equation}
    \tanh x = \frac{e^x - e^{-x}}{e^x + e^{-x}}
\end{equation}
Consideremos la dinámica de activación de una neurona recurrente con decaimiento pasivo lineal (tipo \emph{leaky integrator}):
\begin{equation}
    \dot{x} = -x + \beta \tanh x
    \label{eq:neurona_ode}
\end{equation}
donde:
\begin{itemize}
    \item $x(t) \in \mathbb{R}$ representa el potencial o nivel de activación de la neurona.
    \item $-x$ modela la fuga pasiva hacia el potencial de reposo en ausencia de estímulo.
    \item $\beta > 0$ es el \textbf{peso sináptico recurrente} o ganancia de auto-excitación.
    \item $\tanh x$ es la función de activación no lineal sigmoidal que satura en $\pm 1$ para valores extremos de excitación.
\end{itemize}

Notemos que el sistema posee simetría impar estricta ante reflexiones: si $f(x) = -x + \beta\tanh x$, entonces $f(-x) = -(-x) + \beta\tanh(-x) = x - \beta\tanh x = -f(x)$. Por tanto, el campo es equivariante.

\subsection{Análisis geométrico: intersección de $x$ y $\beta\tanh x$}

Los puntos fijos del sistema satisfacen la condición de equilibrio:
\begin{equation}
    -x + \beta \tanh x = 0 \iff x^* = \beta \tanh(x^*)
    \label{eq:interseccion_neurona}
\end{equation}
Geométricamente, los equilibrios corresponden a los puntos de intersección entre la recta identidad $y = x$ y la función sigmoidal escalada $y = \beta \tanh x$.

Dado que $\tanh(0) = 0$, el origen $x^* = 0$ es siempre una solución para cualquier valor de $\beta$. Para saber si existen soluciones adicionales no nulas, comparamos las pendientes de ambas curvas en el origen:
\begin{equation}
    \left.\frac{d}{dx}(\beta \tanh x)\right|_{x=0} = \beta \operatorname{sech}^2(0) = \beta (1) = \beta
\end{equation}

\begin{figure}[htbp]
\centering
\begin{tikzpicture}[>=Stealth, scale=0.82]
    % Panel a: beta < 1 (beta = 0.6)
    \begin{scope}[shift={(-6.2, 0)}]
        \draw[->, thick, black!60] (-2.2,0) -- (2.2,0) node[right, font=\tiny] {$x$};
        \draw[->, thick, black!60] (0,-1.8) -- (0,1.8) node[above, font=\tiny] {$y$};
        \draw[dashed, black!50, thick] (-1.5,-1.5) -- (1.5,1.5) node[right, font=\tiny] {$y=x$};
        \draw[very thick, black!85, domain=-1.8:1.8, samples=50] plot (\x, {0.6*tanh(\x)});
        \filldraw[black] (0,0) circle (2pt);
        \node[font=\footnotesize, align=center] at (0,-2.1) {\textbf{(a)} $\beta < 1$ ($\beta=0.6$)\\1 equilibrio estable};
    \end{scope}

    % Panel b: beta = 1
    \begin{scope}[shift={(0, 0)}]
        \draw[->, thick, black!60] (-2.2,0) -- (2.2,0) node[right, font=\tiny] {$x$};
        \draw[->, thick, black!60] (0,-1.8) -- (0,1.8) node[above, font=\tiny] {$y$};
        \draw[dashed, black!50, thick] (-1.5,-1.5) -- (1.5,1.5) node[right, font=\tiny] {$y=x$};
        \draw[very thick, black!85, domain=-1.8:1.8, samples=50] plot (\x, {1.0*tanh(\x)});
        \filldraw[black] (0,0) circle (2pt);
        \node[font=\footnotesize, align=center] at (0,-2.1) {\textbf{(b)} $\beta = 1$ (tangencia)\\Bifurcación crítica};
    \end{scope}

    % Panel c: beta > 1 (beta = 1.8)
    \begin{scope}[shift={(6.2, 0)}]
        \draw[->, thick, black!60] (-2.2,0) -- (2.2,0) node[right, font=\tiny] {$x$};
        \draw[->, thick, black!60] (0,-1.8) -- (0,1.8) node[above, font=\tiny] {$y$};
        \draw[dashed, black!50, thick] (-1.5,-1.5) -- (1.5,1.5) node[right, font=\tiny] {$y=x$};
        \draw[very thick, black!85, domain=-1.8:1.8, samples=50] plot (\x, {1.8*tanh(\x)});
        \filldraw[black] (-1.4, -1.4) circle (2pt) node[below left, font=\tiny] {$-x^*$};
        \draw[black, fill=white, thick] (0,0) circle (2pt);
        \filldraw[black] (1.4, 1.4) circle (2pt) node[above right, font=\tiny] {$+x^*$};
        \node[font=\footnotesize, align=center] at (0,-2.1) {\textbf{(c)} $\beta > 1$ ($\beta=1.8$)\\3 equilibrios (biestabilidad)};
    \end{scope}
\end{tikzpicture}
\caption{Determinación geométrica de los puntos de equilibrio de la neurona recurrente $\dot{x} = -x + \beta\tanh x$ mediante la intersección de $y=x$ con $y=\beta\tanh x$.}
\label{fig:intersecciones_red_neuronal}
\end{figure}

Surgen tres comportamientos cualitativos según la ganancia $\beta$:
\begin{enumerate}
    \item \textbf{Caso $\beta < 1$ (subcrítico):} La pendiente en el origen es menor que 1. Como $\tanh x$ es estrictamente cóncava para $x > 0$, la curva $y = \beta\tanh x$ permanece por debajo de la recta $y = x$ en todo el semieje positivo. El único punto de intersección es $x^* = 0$.
    La derivada es $f'(0) = -1 + \beta < 0$, por lo que el origen es un \textbf{atractor asintóticamente estable global}.
    \item \textbf{Caso $\beta = 1$ (umbral crítico):} La curva $y = \tanh x$ es exactamente tangente a la identidad en el origen. El origen es el único punto fijo y es estable con decaimiento subexponencial (\emph{critical slowing down}).
    \item \textbf{Caso $\beta > 1$ (supercrítico):} La pendiente en el origen es mayor que 1. La curva $\beta\tanh x$ comienza por encima de $y = x$ y, debido a que satura asintóticamente hacia el valor finito $\beta$, forzosamente debe cruzar a la identidad en un valor $x_+^* > 0$. Por simetría impar, existe una intersección idéntica en $-x_+^* < 0$.
    El origen $x^* = 0$ pierde estabilidad ($f'(0) = \beta - 1 > 0$, inestable), mientras que los dos nuevos estados activos $\pm x_+^*$ son \textbf{atractores estables}.
\end{enumerate}

\subsection{Verificación analítica con serie de Taylor de $\tanh x$}

Para demostrar analíticamente que el sistema experimenta una bifurcación pitchfork supercrítica canónica en $\beta_c = 1$, expandimos $\tanh x$ en serie de Maclaurin cerca de $x = 0$:
\begin{equation}
    \tanh x = x - \frac{x^3}{3} + \frac{2x^5}{15} - \mathcal{O}(x^7)
\end{equation}
Sustituyendo en la ecuación diferencial \eqref{eq:neurona_ode}:
\begin{align}
    \dot{x} &= -x + \beta \left(x - \frac{x^3}{3} + \mathcal{O}(x^5)\right) \nonumber \\
    &= (\beta - 1)x - \frac{\beta}{3}x^3 + \mathcal{O}(x^5)
    \label{eq:taylor_neurona}
\end{align}
Cerca del punto de bifurcación ($\beta \approx 1$), podemos aproximar $\beta/3 \approx 1/3$:
\begin{equation}
    \dot{x} \approx (\beta - 1)x - \frac{1}{3}x^3
\end{equation}
Haciendo el cambio de parámetro $r = \beta - 1$, recuperamos de forma exacta la forma normal supercrítica $\dot{x} = rx - \frac{1}{3}x^3$. Las ramas de equilibrio no nulas para $\beta > 1$ vienen dadas localmente por:
\begin{equation}
    (\beta - 1)x^* - \frac{1}{3}(x^*)^3 \approx 0 \implies x_\pm^* \approx \pm\sqrt{3(\beta - 1)}
    \label{eq:ramas_analiticas_neurona}
\end{equation}

\begin{figure}[htbp]
\centering
\begin{tikzpicture}[>=Stealth, scale=0.85]
    \draw[->, thick, black!60] (0,0) -- (4.5,0) node[right, font=\small] {$\beta$ (ganancia)};
    \draw[->, thick, black!60] (0,-2.2) -- (0,2.2) node[above, font=\small] {$x^*$};
    
    % Rama trivial
    \draw[ultra thick, black!85] (0,0) -- (1.5,0);
    \draw[very thick, dashed, black!75] (1.5,0) -- (4.0,0);
    
    % Ramas de la pitchfork x* = \pm \sqrt{3(\beta - 1)} => \beta = 1 + (x*)^2 / 3
    \draw[ultra thick, black!85, domain=0:1.6, samples=50] plot ({1.5 + \x*\x*0.9}, \x);
    \draw[ultra thick, black!85, domain=0:1.6, samples=50] plot ({1.5 + \x*\x*0.9}, {-\x});
    
    \filldraw[black] (1.5,0) circle (2.5pt) node[below=4pt, font=\footnotesize] {$\beta_c = 1$};
    \node[above, font=\footnotesize] at (3.2, 1.4) {$x_+^*$ (memoria $+1$)};
    \node[below, font=\footnotesize] at (3.2, -1.4) {$x_-^*$ (memoria $-1$)};
    \node[above, font=\footnotesize] at (0.8, 0.1) {Reposo inactivo};
\end{tikzpicture}
\caption{Diagrama de bifurcación pitchfork de la neurona recurrente en función del parámetro de ganancia sináptica $\beta$. Para $\beta > 1$, la red almacena un bit de memoria binaria biestable.}
\label{fig:bifurcacion_red_neuronal}
\end{figure}

Esta bifurcación es el mecanismo fundamental mediante el cual las redes recurrentes implementan \textbf{memoria asociativa y toma de decisiones biestable}: para $\beta > 1$, el sistema puede conmutar entre los atractores $+x^*$ (estado binario $1$) y $-x^*$ (estado binario $0$).
